\documentclass[11pt,a4paper]{article}
\usepackage[T1]{fontenc}
\usepackage[utf8]{inputenc}
\usepackage{amsmath,amssymb,amsthm}
\usepackage[margin=1in]{geometry}
\usepackage[colorlinks=true,linkcolor=blue,urlcolor=blue]{hyperref}
\theoremstyle{plain}
\newtheorem{theorem}{Theorem}
\newtheorem{lemma}[theorem]{Lemma}
\newtheorem{proposition}[theorem]{Proposition}
\newtheorem{corollary}[theorem]{Corollary}
\theoremstyle{definition}
\newtheorem{remark}[theorem]{Remark}

\title{Recovering the unwritten recurrences of the table OEIS A253375}
\author{Adrian Perez Fontelles\\ \small Independent researcher}
\date{2 September 2026}
\begin{document}
\maketitle

\begin{abstract}
OEIS A253375 is a two-dimensional table $T(n,k)$ whose empirical recurrences are, for several of
its lines, given only as an ORDER --- ``k=2: [order 45]'' and the like --- with the
coefficients in a linked file rather than in the entry. Those conjectures can still be settled,
and the recurrences recovered. A line of the table is a fixed-width or fixed-height array
count, hence a walk count on finitely many vertices, hence satisfies some monic recurrence of
order at most that number; Berlekamp--Massey on twice as many exact terms returns the MINIMAL
one. Where its order equals the order the entry states --- and where the same holds after the
first few terms are dropped --- any recurrence of that order the line satisfies has a
characteristic polynomial that is a multiple of the minimal one and of equal degree, hence is
that polynomial. This note settles column 2, row 2.
\end{abstract}

\noindent\small 2020 Mathematics Subject Classification. 05A15, 11B37, 15A18.\normalsize

\section{The table and the unwritten conjectures}

OEIS A253375 is ``T(n,k)=Number of (n+1)X(k+1) 0..2 arrays with every 2X2 subblock sum nondecreasing horizontally, vertically and antidiagonally ne-to-sw.''. It is read by antidiagonals and begins
\[
81, 450, 450, 2500, 3111, 2500, 11100, 22877,\ \dots
\]
The entry carries, among its empirical claims:
\begin{quote}\small
k=2: [order 45]\\
n=2: [order 79]
\end{quote}
As of the ``Last modified'' line on the live entry (Jul 23 2025 14:01:08, revision 6) these are still
recorded as empirical, and the recurrences themselves are not in the entry text.

\section{Why a line of the table is a walk count}

Fix the index that the line fixes. The entry defines $T(n,k)$ as a number of arrays of a shape
determined by $n$ and $k$, subject to a condition local to a bounded window of consecutive
lines. Substituting the fixed index into the entry's own wording turns the definition into an
ordinary one-parameter array count, and for such a count the admissible windows are the
vertices of a finite digraph and an array is a walk. So the line is a walk count on $S$
vertices, and by the Cayley--Hamilton theorem it satisfies a monic integer recurrence of order
at most $S$.

\section{Recovering the recurrences}

\begin{lemma}\label{lem:bm}
A sequence known to satisfy some linear recurrence of order at most $S$ is determined, with its
minimal recurrence, by its first $2S$ terms; Berlekamp--Massey applied to them returns that
minimal recurrence exactly.
\end{lemma}

\subsection*{Column $k=2$}
Substituting the fixed index into the entry's wording gives an ordinary one-parameter array
count, a walk on $S=450$ vertices. Berlekamp--Massey on $2S=900$ exact terms
returns a minimal recurrence of order $45$ --- the order the entry states --- with
integer coefficients
\[
(c_1,\dots,c_{45})=(9,\ 18,\ -390,\ 363,\ 7557,\ -16924,\ -85080,\ 289038,\ 592626,\ -3022596,\ -2302692,\ 22010974,\ 263586,\ -118108608,\ 60625464,\ 480363315,\ -453795339,\ -1497838078,\ 2041953882,\ 3566209239,\ -6619257735,\ -6304597140,\ 16357594512,\ 7550240128,\ -31527338352,\ -3791317152,\ 47758481952,\ -6659648352,\ -56728278048,\ 20129722880,\ 52207707264,\ -28856183808,\ -36302235648,\ 27669476352,\ 18144427008,\ -18828638464,\ -5799153408,\ 9055982592,\ 719677440,\ -2944438272,\ 244629504,\ 582598656,\ -119439360,\ -53084160,\ 15925248).
\]
so that $a(m)=\sum_{i=1}^{45}c_i\,a(m-i)$ for every $m$. Removing the first
$45+4$ terms and repeating gives order $45$ again, which is what the
identification below needs.

\subsection*{Row $n=2$}
Substituting the fixed index into the entry's wording gives an ordinary one-parameter array
count, a walk on $S=450$ vertices. Berlekamp--Massey on $2S=900$ exact terms
returns a minimal recurrence of order $79$ --- the order the entry states --- with
integer coefficients
\[
(c_1,\dots,c_{79})=(9,\ 65,\ -813,\ -1503,\ 35067,\ -2049,\ -960243,\ 1086240,\ 18711228,\ -35736696,\ -275381736,\ 714658500,\ 3165994812,\ -10469468628,\ -28924095228,\ 120561637986,\ 210319415742,\ -1132860637250,\ -1190202257382,\ 8890114025950,\ 4801619959530,\ -59186272390110,\ -8367164615610,\ 338018085156660,\ -63026938123020,\ -1669263148180740,\ 783202925756340,\ 7168284121132980,\ -5196566728384740,\ -26866295989193748,\ 26010783052803972,\ 88048679140737015,\ -106748031772881279,\ -252329474348288919,\ 371635847674498251,\ 630909828966841993,\ -1116759425621881869,\ -1368846121844868905,\ 2925275798772753429,\ 2549493296941549716,\ -6718378143067818744,\ -3990279397586265900,\ 13573200380588028348,\ 5003881995007001088,\ -24158694316978174512,\ -4356492449072602896,\ 37885311670133837328,\ 702180134840233920,\ -52281979345965347328,\ 6495657021495425856,\ 63334702238132746176,\ -16079014545805946880,\ -67090134345597242112,\ 25187434137260649728,\ 61801025531016991488,\ -30523610020098282496,\ -49124893907374282752,\ 30216827194676075520,\ 33334146513723036672,\ -24903026047945359360,\ -19010417886235054080,\ 17168181780647301120,\ 8895860007639920640,\ -9868786556970516480,\ -3276824012368773120,\ 4685347765788917760,\ 868947223463608320,\ -1807894348725288960,\ -120651665681940480,\ 553134494117265408,\ -17459153936842752,\ -129178265777602560,\ 14777776041099264,\ 21637807214690304,\ -4007539482034176,\ -2315214083063808,\ 570630428688384,\ 118881339310080,\ -35664401793024).
\]
so that $a(m)=\sum_{i=1}^{79}c_i\,a(m-i)$ for every $m$. Removing the first
$79+4$ terms and repeating gives order $79$ again, which is what the
identification below needs.

\begin{theorem}
For each line above, the displayed recurrence holds for every index, and it is the recurrence
the entry records.
\end{theorem}

\begin{proof}
It holds by Lemma~\ref{lem:bm}. For the identification, the entry asserts that the line
satisfies SOME recurrence of the stated order, possibly only from some index on. Let $q$ be its
characteristic polynomial and $q_{\min}$ the minimal polynomial of the tail. Then
$q_{\min}\mid q$, and the second Berlekamp--Massey run --- on the line with its first few
terms removed --- shows $\deg q_{\min}$ equals the stated order, which is $\deg q$. Both are
monic, so $q=q_{\min}$.
\end{proof}

\section{Verification}

Three checks.

First, each line's model was compared against the entry's own published values for that line,
taken from the antidiagonal DATA read in either orientation and from the ``Table starts''
block, and agrees at every available term. This is the only point at which reading the entry's
English enters, and a misreading gives different counts.

Second, each recovered recurrence was evaluated directly on the model's values, with no
Berlekamp--Massey involved, and holds at every index tested.

Third, each order was computed twice by independent routes --- modulo a $61$-bit prime, where
every intermediate is a machine word, and again in exact rational arithmetic --- and once more
on the line with its first few terms dropped, which is what the identification requires. All
agree.

\begin{thebibliography}{9}
\bibitem{oeis} The OEIS Foundation, \emph{The On-Line Encyclopedia of Integer
Sequences}, \url{https://oeis.org}, 2026. Sequence A253375.
\bibitem{massey} J.~L.~Massey, Shift-register synthesis and BCH decoding, \emph{IEEE
Trans. Inform. Theory} \textbf{15} (1969), 122--127.
\bibitem{stanley} R.~P.~Stanley, \emph{Enumerative Combinatorics}, Volume 1, 2nd ed.,
Cambridge University Press, 2011. (Transfer-matrix method, Section 4.7.)
\bibitem{horn} R.~A.~Horn and C.~R.~Johnson, \emph{Matrix Analysis}, 2nd ed., Cambridge
University Press, 2013.
\end{thebibliography}

\end{document}
